%%%%%%%%%%%%%%%%%%%%%%%%%%%%%%%%%%%%%%%%%%%%%%%%%%%%%%%%%%%%%%%%%%%%%%%%%%%%%%%%
%%
\cleardoubleoddpage%  Make sure to start each chapter on a new odd page
\chapter{Introduction}
\label{ch:introduction}


\section{Clinical Motivation}
\label{sec:introduction:motivation}

% TODO: cite Wong2014 (Lancet Glob Health) once added to references.bib --
% canonical reference for AMD as the leading cause of irreversible central
% vision loss in industrialised countries; currently the sentence rests
% indirectly on Flaxman2020.
Age-related macular degeneration (AMD) is the leading cause of irreversible
central vision loss in older adults in industrialised countries and ranks as
the third most frequent cause of moderate-to-severe vision impairment and the
fourth most frequent cause of blindness worldwide
\citep{Flaxman2020}. Roughly 200~million people are
estimated to live with some form of AMD today, with projections reaching
288~million by 2050 \citep{Boopathiraj2024,Trincao2024}. Among these, the
late dry stage --- termed \emph{Geographic Atrophy} (GA) --- is the dominant
cause of severe visual decline in the absence of choroidal neovascularisation
\citep{Vallino2024}. Approximately one million people in the United States and
five million worldwide are affected by GA, with around 160{,}000 new diagnoses
per year in the United States alone \citep{Singh2025,Vallino2024}. Crucially,
AMD is not classified as an avoidable cause of vision loss in the WHO/GBD
framework \citep{Flaxman2020}; the principal clinical lever therefore lies in
\emph{managing} progression rather than in primary prevention.

The disease burden is concentrated in the older population. Prevalence of GA
rises by roughly a factor of four per decade after age fifty
\citep{Singh2025}, and population-based screening data show overall AMD
prevalence increasing from 3.6\,\% in the fifth decade of life to 34.0\,\%
above the age of eighty \citep{Song2025}. GA is bilateral in approximately
39\,\% to 50\,\% of patients at the time of initial diagnosis, with strong
inter-eye correlation \citep{Yehoshua2011,Boyer2017}. The combination of an
aging population, high bilaterality, and irreversible foveal involvement
defines GA as a long-horizon, individually progressive disease for which the
clinically relevant question is not \emph{whether} a given patient will
deteriorate but \emph{how fast} and \emph{where} in the macula the lesion
will expand.

For decades, the therapeutic landscape for GA was empty: a series of phase
1--3 trials targeting neuroprotection, the visual cycle, antiamyloid pathways,
and cell replacement consistently failed to slow lesion growth
\citep{Lad2023}. This changed in 2023 with the U.S. Food and Drug
Administration's approval of two complement-pathway inhibitors,
pegcetacoplan (anti-C3) and avacincaptad pegol (anti-C5)
\citep{Lad2023,Singh2025,Trincao2024}. Both drugs reduce the rate of GA
enlargement by roughly 14\,\% to 27\,\% over twelve to twenty-four months but
neither arrests nor reverses the underlying atrophy, and best-corrected
visual acuity remains unchanged at the trial endpoints
\citep{Lad2023,Trincao2024}. The clinical implication is twofold. First,
treatment timing now matters: an intervention that only slows progression is
most valuable when administered before foveal involvement
\citep{Boyer2017,Vallino2024}. Second, intravitreal injections carry both
patient burden and procedural risk, so identifying which eyes will progress
rapidly enough to justify therapy is the critical decision support task. As
\citet{Vallino2024} explicitly note in their conclusions, the integration of
artificial intelligence with structural OCT biomarkers offers the most
promising avenue toward predictive probabilities of GA development over time.
This thesis pursues that avenue.


\section{Problem Statement}
\label{sec:introduction:problem}

Predicting GA progression from imaging is naturally framed as a
\emph{spatiotemporal forecasting} task on the en-face projection of optical
coherence tomography (OCT) volumes. For each eye, longitudinal imaging yields
a sequence of macula-centred grids on which the lesion outline and the
underlying retinal layer geometry are traced; given the most recent observed
state and a target horizon $\Delta t$, the goal is to predict the full
spatial configuration of the lesion at the future visit. This per-eye,
spatially resolved framing is distinct from the cohort-level conversion-time
prediction common in prior AMD deep-learning work \citep{Vogl2021}: the
deliverable is not a single risk score but the entire next image of the
disease.
% TODO: cite SchmidtErfurth2018 (IOVS, "Prediction of individual disease
% conversion in early AMD using artificial intelligence") and
% Bogunovic2017 (IOVS, "Machine learning of the progression of intermediate
% AMD based on OCT imaging") once added to references.bib -- both are
% direct AI-on-OCT precedents for the cohort-level conversion framing,
% Bogunovic is the supervisor's own institutional precedent.

The clinical reality of OCT follow-up imposes three properties on the data
that the model must respect. First, follow-up visits are
\emph{irregularly spaced} in time and the spacing varies substantially
between patients and between consecutive visits of the same patient. There
is no shared temporal origin across eyes, which rules out methods that
assume a fixed sampling rate or a common $t = 0$. Second, the underlying
disease state is naturally \emph{multi-channel}: the binary lesion mask
captures the cRORA (complete retinal pigment epithelium and outer retinal
atrophy) region itself, while the surrounding ten retinal layer boundaries
encode the structural cascade in which the atrophy is embedded
\citep{Chu2022,Vallino2024}. Both information sources are required because
prognostic biomarkers documented in the clinical literature, such as
RPE--Bruch's-membrane thickening and outer nuclear layer thinning, are not
contained in the lesion mask alone. Third, every patient carries
\emph{covariates} --- in particular age and sex --- that modulate baseline
risk and progression speed and that should enter the dynamics on equal
footing with the imaging signal.

The combination of these requirements --- a per-eye spatial state forecast,
under arbitrary $\Delta t$, on a vector-valued field, conditioned on patient
context --- is exactly the kind of task for which conventional segmentation
networks are ill-suited.
% TODO: cite Ronneberger2015 (MICCAI, "U-Net: Convolutional Networks for
% Biomedical Image Segmentation") once added to references.bib -- the
% U-Net mention below should carry an inline citation.
A standard U-Net or similar dense predictor
provides no principled mechanism for ingesting an elapsed-time scalar, no
notion of the local-dynamics structure that drives lesion-boundary growth,
and no separation between the static disease identity of the patient and
the time-evolving state. The next section argues that the problem instead
fits naturally into the framework of neural partial differential equation
(PDE) solvers.


\section{Why Neural PDE Solvers?}
\label{sec:introduction:why-pde}

GA expansion is, to a first approximation, a process of \emph{local boundary
dynamics}. \citet{Yehoshua2011} demonstrated empirically that the radius of
a GA lesion grows at a roughly constant rate while its enclosed area scales
quadratically: writing $A(t)$ for the lesion area and $r(t)$ for an
effective radius,
\begin{equation}
  \frac{\mathrm{d}r}{\mathrm{d}t} \approx \mathrm{const.}
  \quad\Longleftrightarrow\quad
  \sqrt{A(t)/\pi} \;=\; r_0 + v\, t,
  \label{eq:intro:boundary-growth}
\end{equation}
which holds within the clinical follow-up horizon of six months to slightly
over two years and has been re-used in subsequent biomarker work
\citep{Chu2022}. The interior of the lesion is already atrophic and the
exterior is healthy; both are quasi-static, and all dynamics are
concentrated at the active boundary. This is the structural fingerprint of
a system governed by a temporal partial differential equation with strong
spatial localisation.

The same boundary, however, does not advance uniformly in all directions.
GA growth is \emph{anisotropic}: lesions progress faster towards the retinal
periphery than towards the fovea, peak progression rate occurs at
approximately one millimetre from the foveal centre when the lesion grows
foveally, and complement inhibition has a spatially asymmetric treatment
effect \citep{Singh2025}. The forward operator that maps the present state
to the future state therefore has both \emph{local} and
\emph{spatially adaptive} character --- two inductive biases that classical
finite-difference and finite-volume PDE solvers exploit explicitly through
local stencils and adaptive meshes, and that generic dense convolutional
predictors must learn implicitly from data.

% TODO: in the autoregressive-vs-operator contrast below, the operator
% family (FNO, DeepONet) is intentionally deferred to Chapter 2.3
% per THESIS_STRUCTURE.md; once Li2021 (FNO) and Lu2021 (DeepONet) are
% added to references.bib, decide whether a one-sentence acknowledgement
% of the operator alternative belongs here as well.
Two recent neural PDE solver families directly encode these biases.
\citet{Brandstetter2022} introduced \emph{message-passing neural PDE
solvers} (MP-PDE), in which a graph neural network performs an
% TODO: cite Battaglia2018 (relational inductive biases) and
% SanchezGonzalez2020 (Learning to Simulate) once added to references.bib --
% the encode--process--decode pattern predates Brandstetter2022 and is
% conventionally attributed to those works.
encode--process--decode rollout that produces a residual Euler update of the
form $u_{t+\Delta t} = u_t + \Delta t \cdot f_{\theta}(u_t, \Delta t)$. The
message-passing layers act as learned, non-linear analogues of classical
local stencils, and the temporally conditioned residual form provides a
natural place to inject elapsed time into the dynamics.
\citet{Hu2024} extended this idea with the \emph{moving-mesh PDE solver}
(MM-PDE), which couples a uniform-mesh GNN with a second GNN operating on a
dynamically moved mesh produced by a data-free mesh mover; the moving mesh
concentrates discretisation density in regions of large solution gradient,
which on GA imagery is precisely the active lesion boundary. Adapting these
two frameworks to longitudinal OCT bridges the methodological PDE-solver
literature with the applied medical-imaging deep-learning literature
\citep{Vogl2021,Vallino2024}, where prior work has remained anchored in
cohort-level statistical atlases and time-to-conversion survival models
rather than in per-eye autoregressive forward operators on the imaging
state itself.

\paragraph{Research question.}
\label{par:introduction:research-question}
This thesis is organised around a single question:

\begin{quote}
  \itshape Given the lack of regenerative therapies for Geographic
  Atrophy and the recent approval of complement inhibitors that only
  slow lesion expansion, can the adaptation of neural PDE solvers to the
  irregular intervals of longitudinal OCT yield per-eye spatial
  forecasts accurate enough to inform patient stratification and
  treatment timing?
\end{quote}

The remaining chapters decompose this question into a sequence of
methodological adaptations and empirical evaluations.


\section{Contributions}
\label{sec:introduction:contributions}

The work presented in this thesis makes the following contributions:

\begin{itemize}
  \item \textbf{First application of moving-mesh neural PDE solvers to
    medical imaging.} The MM-PDE framework of \citet{Hu2024} is
    transferred from synthetic physical benchmarks to longitudinal OCT
    imaging of Geographic Atrophy, and is adapted to the constraints of
    a clinical cohort drawn from the Medical University of Vienna.

  \item \textbf{$\Delta t$-conditioned residual formulation for irregular
    visit schedules.} The fixed-step rollout of the original MP-PDE and
    MM-PDE architectures is replaced with a per-window
    $\Delta t$-conditioned Euler residual,
    \begin{equation}
      u_{t+\Delta t} \;=\; u_t \;+\; \Delta t \cdot f_{\theta}(u_t, \Delta t),
      \label{eq:intro:dt-residual}
    \end{equation}
    which removes the assumption of a shared temporal origin across
    patients and supports inference at clinically arbitrary horizons.

  \item \textbf{Multi-channel generalisation of the MM-PDE pipeline.}
    The moving-mesh and interpolation components of MM-PDE are extended
    from a scalar field to the eleven-channel disease state (binary
    lesion mask plus ten retinal layer boundaries) by reformulating the
    monitor function in terms of the Frobenius norm of the per-channel
    spatial gradient.

  \item \textbf{Patient covariate conditioning.} Age and sex are
    incorporated as graph-level attributes that are broadcast per node
    into both the embedding and every message-passing layer, so that
    patient context modulates the dynamics on equal footing with the
    learned per-step solver state.

  \item \textbf{Surrogate equation encoder.} A lightweight convolutional
    encoder compresses the ten retinal layer maps into a compact global
    embedding that is concatenated to the per-node feature vector; this
    embedding plays the role of the equation-feature vector
    $\theta_{\mathrm{PDE}}$ used in MP-PDE \citep{Brandstetter2022} and
    provides the GNN with a learned surrogate for the otherwise unknown
    PDE coefficients of GA dynamics.

  \item \textbf{Empirical validation on a clinical longitudinal OCT
    cohort.} A persistence baseline, a single-branch GNN (MP-PDE style),
    and the dual-branch MM-PDE adaptation are compared head-to-head with
    matched parameter budgets on per-pixel error, on lesion-mask Dice
    and IoU at a clinically meaningful horizon, and on a battery of
    component ablations covering each of the contributions above.
\end{itemize}


\section{Thesis Outline}
\label{sec:introduction:outline}

The remainder of this thesis is organised as follows. Chapter~\ref{ch:background}
provides the necessary background, surveying GA pathology and OCT imaging,
the relevant subset of numerical and neural PDE solver theory, and the two
specific architectures (MP-PDE, MM-PDE) on which the present work builds.
Chapter~\ref{ch:data} describes the Medical University of Vienna GA cohort,
the eleven-channel state representation, and the spatial, temporal, and
demographic preprocessing pipeline.
Chapter~\ref{ch:method} details the proposed adaptation: the
$\Delta t$-conditioned residual operator, the multi-channel pipeline, the
data-free mesh mover for vector-valued GA states, the dual-branch
correction-style architecture, the patient covariate conditioning, the
surrogate equation encoder, and the training protocol.
Chapter~\ref{ch:experiments} reports the empirical evaluation, comprising
baseline comparisons, ablations of each contribution, moving-mesh quality
diagnostics, qualitative rollout analyses, and computational-cost
measurements.
Chapter~\ref{ch:discussion} interprets the results in clinical and
methodological terms and articulates the principal limitations.
Chapter~\ref{ch:conclusion} summarises the contributions and outlines
directions for future work.

% TODO: confirm chapter labels (\ref{ch:background} etc.) match the
% \label commands defined in the chapter source files once they are written;
% currently these are placeholders consistent with the planned structure.


%%
%%%%%%%%%%%%%%%%%%%%%%%%%%%%%%%%%%%%%%%%%%%%%%%%%%%%%%%%%%%%%%%%%%%%%%%%%%%%%%%%
